\section{a) Runge-Kutta method of $4^{th}$ order and Adams PC}
\subsection{Problem}
We are given following equations:
\[ \frac{dx_1}{dt} = x_2 + x_1(0.5 - x_1^2 - x_2^2) \]
\[ \frac{dx_2}{dt} = -x_1 + x_2(0.5 - x_1^2 - x_2^2) \]

And we have to determine the trajectory of the motion on interval $[0, 15]$ with following initial conditions:
$ x_1(0) = 8; x_2(0) = 9 $
In this section we will use Runge-Kutta method of $4^{th}$ order and Adams PC with different step-sizes until we find an optimal constant step size - when the decrease of the step size does not influence the solution significantly.

\subsection{Theoretical Introdution}
Determining the trajectory of the motion in this example can be mathematically described as solving a system of first order ordinary differential equations with given initical conditions. That is why we use Runge-Kutta method of order 4 which deals with exactly this kind of problem.
\subsubsection{Runge-Kutta method of order 4}
Any single step method are defined by the following formula, where h is a fixed step-size:
\[ y_{n+1} = y_n + \Phi_f(x_n, y_n, h) \]
where:
\[ x_n = x_0 + nh \]
\[ y(x_0) = y_0 = y_a \]
where $y_a$ is given.
Runge-Kutta method of order 4 - RK4 method looks like this:
\[ y_{n+1} = y_n + \frac{1}{6}(k_1 + 2k_2+2k_3+k_4) \]
\[ k_1 = f(x_n, y_n) \]
\[ k_3 = f(x_n + \frac{1}{2}h, y_n + \frac{1}{2}hk_1) \]
\[ k_2 = f(x_n + \frac{1}{2}h, y_n + \frac{1}{2}hk_2) \]
\[ k_4 = f(x_n + h, y_n + hk_3) \]


\subsection{Adams PC method}
Adams method with $P_kEC_kE$ algorithm has the following form: \\
P: \[  y_n^{[0]} = y_{n-1} + h\sum_{j=1}^k \beta_jf_{n-j} \]
E: \[  f_n^{[0]} = f(x_n, y_n^{[0]}) \]
C: \[  y_n = y_{n-1} + h \sum_{j=1}^k \beta_j^* f_{n-j} + h \beta_0^* f_n^{[0]}\]
E: \[  f_n = f(x_n, y_n) \]

When compared to explicit and implicit adams method PC metho has smaller absolute stability intervals than implicit method, while having significant advantage in calculations cost.



